% -*- coding: utf-8 -*-

\begin{zhixie}
\begin{spacing}{1.5}
本文的完成离不开许多老师、同学与家人的支持与帮助。首先，我要衷心感谢指导教师在选题确定、方案设计、系统实现和论文写作全过程中给予的耐心指导与严格要求。老师在长期记忆智能体评测框架的研究思路、工程实现边界以及论文结构组织方面提出了大量具体而关键的建议，使我能够在复杂的工程实践中不断修正方向、完善细节。

其次，感谢课题组与同学们在项目讨论、系统调试和文档整理过程中给予的帮助。围绕 O-Mem、MemBox 适配、实验脚本重构、结果分析与环境排障的多次交流，使我逐步形成了对长期记忆评测问题更系统的认识，也帮助我在遇到复杂工程问题时保持耐心和严谨。

同时，我也感谢家人在本科阶段始终给予的理解、支持与鼓励。正是这种稳定而温暖的支持，使我能够在论文撰写和系统实现过程中保持持续投入，并顺利完成本阶段学习任务。

最后，向所有在百忙之中评阅本文并提出宝贵意见的老师和专家致以诚挚的感谢。
\end{spacing}
\end{zhixie}
